\section{Unbounded resources: antichains over a well-quasi-order}
\label{sec:wsts}

Four contributing efforts arrive at the same mathematical shape, and the shape
is worth stating before any of them. A transition system is \emph{monotone} if
having more resources never disables a step: $s\preceq s'$ and $s\to s'$'s
source can fire implies the larger state can fire the same transition and land
at least as high. If in addition $\preceq$ is a well-quasi-order --- every
infinite sequence contains $i<j$ with $s_i\preceq s_j$ --- then the set of
states from which something bad is reachable is upward closed, and an upward
closed set in a well-quasi-order has a \emph{finite} basis.

That is the whole idea. The unsafe region of an infinite-state system is
described exactly by finitely many minimal states, and a backward search
terminates because it cannot admit an infinite antichain.

\begin{lemma}[Dickson]\label{lem:dickson}
Every infinite sequence in $\N^d$ contains indices $i<j$ with $x_i\le x_j$
coordinatewise. In particular every antichain in $\N^d$ is finite.
\end{lemma}
\begin{proof}
Every infinite sequence of naturals has an infinite nondecreasing subsequence:
if some value recurs infinitely often take the constant subsequence, and
otherwise choose successively larger values from later positions. Apply this to
the first coordinate, then to the second coordinate of the retained
subsequence, and so on. After $d$ restrictions all coordinates are
nondecreasing, and any two positions of the final subsequence give the pair.
The case $d=0$ is immediate.
\end{proof}

Only one of the four proves this; the others cite it. It is stated here once.
The same is true of the predecessor computation, which every one of them needs
and every one of them derives independently.

\begin{lemma}[Minimal predecessor]\label{lem:backpre}
For a transition consuming $c$ and producing $p$, let
$\pre_t(b)=c+(b-p)^{+}$ where $(\cdot)^{+}$ is coordinatewise positive part.
Then $\{m: \exists m'\ge b,\ m\xrightarrow{t}m'\}=\Up{\{\pre_t(b)\}}$, and
firing $t$ from $\pre_t(b)$ lands at $\max(b,p)\ge b$.
\end{lemma}
\begin{proof}
The condition is $m\ge c$ together with $m-c+p\ge b$. Coordinatewise: where
$b_i\le p_i$ the enabling requirement $m_i\ge c_i$ already suffices; where
$b_i>p_i$ the two inequalities together say $m_i\ge c_i+b_i-p_i$. That is the
displayed maximum. At the threshold the endpoint is $(b-p)^{+}+p=\max(b,p)$.
\end{proof}

The four efforts write this four ways --- as a maximum, as a positive part, as
$\max(a,a+b-z)$, and abstractly through a \emph{strong upward compatibility}
axiom --- and they are the same function.

\subsection{One initial state, or all of them}

Here the four separate, and the distinction is the most consequential thing in
this section because it is invisible in the certificate's shape.

Three of the efforts compute the unsafe region for \emph{every} initial state.
The fourth binds its receipt to one, by adding a third obligation: every basis
element must fail to be below the given $m_0$. Its certificate is otherwise
identical in form --- the same basis, the same target coverage, the same
backward closure --- and it does not require the basis to be minimal or even an
antichain, because for its purpose a redundant over-approximation is enough.

That effort then says exactly what it has left on the table:

\begin{quote}\small
There is a reusable theorem hiding in this certificate. Once
[target coverage] and [backward closure] are checked, the same $A$ proves
safety for \emph{every} initial marking outside $\up A$. The prototype binds
each receipt to one initial marking, but a Lean port should expose the
stronger invariant lemma and then specialize it.
\end{quote}

Another effort in the same round delivers precisely that theorem, checked, with
the extra ingredient the stronger statement needs. The merge states it once, in
the strong form:

\begin{theorem}[Exact frontier]\label{thm:wsts-exact}
Let $B$ be finite with: every bad target above some $b\in B$; for every $b\in B$
and transition $t$, some $a\in B$ with $a\le\pre_t(b)$; and for every $b\in B$
a recorded run $u_b$ enabled at $b$ whose endpoint covers an original target.
Then for every state $m$,
\[
 \text{bad is reachable from } m \iff \exists b\in B,\ b\le m .
\]
If $B$ is an antichain its members are exactly the minimal unsafe states.
\end{theorem}
\begin{proof}
Right to left: choose $b\le m$ and run $u_b$; monotonicity lifts the run from
$b$ to $m$ and the endpoint still covers a target. Left to right: induct
backwards along a bad execution. Its last state is in $\Up B$ by target
coverage; if a successor covers $b\in B$ via $t$ then by
Lemma~\ref{lem:backpre} the predecessor covers $\pre_t(b)$, and backward
closure supplies $a\in B$ below it. For minimality, if $x\le b\in B$ is unsafe
then exactness gives $a\in B$ with $a\le x\le b$, and incomparability forces
$a=b$.
\end{proof}

\begin{remark}[Two halves, and why neither suffices]
Target coverage and backward closure alone prove $\Up B$ is an
\emph{over-approximation} of the unsafe set --- enough for safety, which is why
the one-initial-state certificate needs nothing more. The recorded runs alone
prove an \emph{under-approximation}. Only together do they give exactness, and
only exactness gives a threshold: the least population at which a system
becomes unsafe is a statement about the minimal basis, not about any
over-approximation of it. A frontier computed under a fixed-initial-state
assumption and then reported as global would erase exactly that threshold, and
one of the efforts warns against this in those words.
\end{remark}

\begin{invariant}[A family-relative certificate is not a frontier]
\label{inv:frontier}
A basis proved safe relative to one initial state, or pruned using an invariant
that holds only there, must not be published as the minimal unsafe basis. The
two claims have different strengths and the certificates are distinguishable
only by the presence of the positive runs.
\end{invariant}

\begin{theorem}[Termination]\label{thm:wsts-terminates}
Without operational cutoffs the backward worklist terminates, and returns
either a basis satisfying Theorem~\ref{thm:wsts-exact} or a concrete bad
execution.
\end{theorem}
\begin{proof}
Each admitted candidate lies outside the upward closure of the current
antichain, and deleting a dominated member does not shrink the represented set
because the new smaller member covers its whole cone. If infinitely many were
admitted, Lemma~\ref{lem:dickson} would put an earlier one below a later one,
whose cone was still represented at that time --- contradicting admission. So
finitely many are admitted, each scheduling finitely many predecessor tasks,
and the queue empties. Every admitted entry carries a genuine run by
Lemma~\ref{lem:backpre} and monotonicity; at termination every surviving entry
has been expanded and every predecessor is covered.
\end{proof}

The proof is qualitative, and the efforts say so: it bounds nothing. A
terminating coverability method can still be hopeless on a given input, and one
of them exhibits the reason directly. A transfer step consuming $(1,0)$ and
producing $(0,1)$ with target $(0,N)$ has minimal unsafe basis
$\{(0,N),(1,N-1),\ldots,(N,0)\}$ --- $N+1$ pairwise incomparable vectors. Any
certificate format that lists the basis explicitly must contain all of them.

\subsection{The same three vectors, three times}

The four efforts were prepared independently, and three of them chose mutual
exclusion as their readable worked example. Their computed bases are:
\[
 \{(0,0,2),(1,1,1),(2,2,0)\},\qquad
 \{(0,2,0),(1,1,1),(2,0,2)\},\qquad
 \{(0,2,0),(1,1,1),(2,0,2)\}.
\]
Under the coordinate bijection matching \emph{waiting} to \emph{idle} and
\emph{permit} to \emph{lock}, the first two are the same net and the same
antichain, vector for vector; the third is a third presentation of it. One
effort derives ``$(n,1,0)$ is safe for every $n$'' from it and another checks
its own version at $n=10^{30}$.

This is one result, and this document records it once. It is also a warning
about the shape of the evidence in this collection: three independent packages
producing an identical three-element antichain is not three confirmations of
anything, and a provenance ledger that counts it three times is wrong by a
factor of three on its most quotable example.

\begin{remark}[A field-name collision that is not a naming quibble]
Two of the efforts use the field names \code{basis}, \code{target\_cover} and
\code{predecessor\_cover} for objects with \emph{different} semantics: in one
the basis is minimal and every element carries a positive run, in the other it
may be redundant and carries none. A merge that unified the schemas on the
strength of matching field names would silently upgrade a family-relative
receipt into a frontier. Appendix~\ref{app:schemas} keeps them apart.
\end{remark}

\subsection{Two ways to survive a large instance}

The efforts diverge again on what to do when an instance is big, and the two
answers are complementary rather than competing.

One accelerates the \emph{search}. It proves a power-summary law for repeating
a transition $n$ times and a finite loop-star bound, then emits a witness
containing a \code{repeat} node. The checker is unchanged --- no new trusted
rule --- and a target of $10^{100}$ is certified by a three-node witness.

The other changes the \emph{certificate language}. Its problem is not the
witness but the local predecessor obligation: a merge transition with $d$
residual counters and threshold $h$ has $\binom{h+d-1}{d-1}$ minimal
predecessors, which at $d=8,h=100$ is $26{,}075{,}972{,}546$. Enumerating them
to discover they are all already covered is wasted work, so it adds a receipt
tag whose leaves point directly at global basis indices. The resulting
certificate is $1{,}063$ bytes.

The cost of having only the second is visible in its own corpus: its
$10^{100}$ fixture needs $334$ nodes and $39{,}194$ bytes, because its doubling
loop is unrolled where the first effort would have written one \code{repeat}.
The cost of having only the first is that it cannot express the dense-predecessor
case at all. A production worker wants both, and they compose --- one compresses
the witness, the other compresses the obligation.

\subsection{How far the model stretches}

The four admit different transition rules, and they form a ladder. The
narrowest is the ordinary net: $m'=m-c_t+p_t$ under $m\ge c_t$, with constant
effect. Adding a finite control is routine. The widest admits
$x'=A(x-a)+b$ with $A$ a nonnegative integer matrix, which buys resets (a zero
row), transfers and merges (several unit entries in one row), and duplication
(an entry above one) --- none of which has constant effect, and the narrowest
effort excludes them by name for exactly that reason.

A different axis is the order itself. One effort leaves $\N^d$ entirely for
finite-control \emph{lossy FIFO} channels, ordered by the subsequence embedding
$u\sqsubseteq w$, which is a well-quasi-order by Higman's lemma rather than
Dickson's. The predecessor law is correspondingly different --- for a receive,
$\Pre(\Up u)=\Up{\{au\}}$; for a send, $\Up{\{v\}}$ where $u=va$ or $u$ itself
--- and it is proved rather than asserted. The abstract monotonicity axiom is
what lets one termination proof cover both domains.

\begin{remark}[Lossy is not reliable]
A counterexample under lossy semantics is not a counterexample to the reliable
system. The recorded fixture is two controls and one receive: from queue $ba$
the lossy system may drop $b$ and then receive $a$, reaching the bad control,
while the reliable system's receive is blocked. Both facts are true, and only
one of them is about the program a user wrote.
\end{remark}

\subsection{Finiteness for a different reason}

One effort adds a worker in this family that does \emph{not} rest on a
well-quasi-order at all, and says so plainly: ``For counters, finiteness comes
from a well quasi-order on natural vectors. For names, it comes from
quotienting by renamings that preserve everything the program can observe.
These are different reasons for finiteness.''

The model is a register machine over an infinite alphabet of atoms, with a
finite control and finitely many registers each holding an atom or nothing.
What is finite is not a basis but the set of \emph{equality patterns} among
registers and constants --- which atoms coincide with which --- and the count is
bounded by a Bell number, $B_{k+c+1}$ for $k$ registers and $c$ constants. Two
configurations with the same pattern are related by a renaming the machine
cannot detect, so equivalence of two deterministic such machines over all
finite words is decidable by exploring patterns.

The implementation detail that matters is a discipline rather than an
algorithm, and it generalises past this worker: the executable representation
numbers atoms as integers, but \emph{no arithmetic or comparison on those
integers is available to the modelled program}. A Python sort used to pick a
representative is a tool of the checker, not an operation of the machine. A
source bridge that forgot this would prove a theorem about a different, more
powerful language.

\subsection{What this section does not claim}

Coverability is not reachability. With one transition $x\mapsto x+2$ from
$x=0$, the target $x\ge1$ is coverable and the exact marking $x=1$ is not
reachable; a source adapter that translates an equality goal into an upward
coverage query has changed the theorem. Nor is any of this a claim about
fairness, deadlock-freedom under a scheduler, recurring visits, or exact-state
reachability --- all of which need different predicates and none of which is
imported by a coverability result.

Two of the four also report defects found in their own harnesses rather than in
their mathematics: a Python equality that let a Boolean compare equal to an
integer, and a reporting key collision that replaced a list of basis vectors by
its length. Both are scoped honestly by their authors --- ``an engineering
correction, not a discovered false Lean proof''; ``a reporting defect, not
evidence of a false region certificate'' --- and both are the kind of thing that
only shows up when someone actually runs the code.

And, as everywhere in this document, none of it has been checked by Lean. Three
of these four ship no Lean at all, two of them stating that as a policy rather
than an accident; the fourth ships a single small specimen whose one-step
complement-invariance lemma is the shared kernel of all four soundness
theorems, and is the right first thing to formalise.
